\documentclass{article}
\usepackage{graphicx}
\usepackage{float}  
\usepackage[utf8]{inputenc}
\usepackage[includeheadfoot, margin=1em,headheight=2em]{geometry}
\usepackage{titling}
\usepackage{hyperref}
\geometry{a4paper, left=2cm, right=2cm, top=2cm, bottom=2cm}
\usepackage{graphicx}
\providecommand{\versionnumber}{0.0.1}
\usepackage{enumitem}
\usepackage{array}
\newcolumntype{P}[1]{>{\centering\arraybackslash}p{#1}}
\renewcommand{\arraystretch}{1.5} % Default value: 1
\setlength{\droptitle}{-6em}
\usepackage{capt-of}
\usepackage{setspace}

%font
\usepackage[defaultfam,tabular,lining]{montserrat}
\usepackage[T1]{fontenc}
\renewcommand*\oldstylenums[1]{{\fontfamily{Montserrat-TOsF}\selectfont #1}}

%custom bold 
\usepackage[outline]{contour}
\usepackage{xcolor}
\newcommand{\custombold}{\contour{black}}

%table colors
\usepackage{color, colortbl}
\definecolor{Blue}{rgb}{0.51,0.68,0.79}
\definecolor{LightBlue}{rgb}{0.82,0.87,0.90}
\definecolor{LighterBlue}{rgb}{0.93,0.95,0.96}

\usepackage{caption}
\captionsetup[figure]{labelformat=empty}

%Header
\usepackage{fancyhdr, xcolor}
\pagestyle{fancy}
\let\oldheadrule\headrule% Copy \headrule into \oldheadrule
\renewcommand{\headrule}{\color{Blue}\oldheadrule}% Add colour to \headrule
\renewcommand{\headrulewidth}{0.2em}
\fancyhead[L]{Manuale Sviluppatore}
\fancyhead[C]{Cybersorceres}
\fancyhead[R]{versione \versionnumber}

\title{\Huge{\textbf{Manuale Sviluppatore}}\vspace{-1em}}
\author{CyberSorcerers Team}
\date{}
\begin{document}
\maketitle
\vspace{-3em}
\begin{figure}[h]
  \centering
  \includegraphics[width=6cm, height=6cm]{documenti/logo rotondo.png}
  \label{fig:immagine}
\end{figure}

\vspace{6em}
\large{
\begin{center}
    \begin{tabular}{P{24em}}
        \rowcolor{Blue}
        \textbf{Membri del team:}\\
        \rowcolor{LighterBlue}
        \custombold{Sabrina Caniato}\\
        \rowcolor{LightBlue}
        \custombold{Giulia Dentone}\\
        \rowcolor{LighterBlue}
        \custombold{Nicola Lazzarin}\\
        \rowcolor{LightBlue}
        \custombold{Giovanni Moretti}\\
        \rowcolor{LighterBlue}
        \custombold{Andrea Rezzi}\\
        \rowcolor{LightBlue}
        \custombold{Samuele Vignotto}\\
    \end{tabular}
\end{center}

\begin{center}
    \begin{tabular}{l c c}
        \rowcolor{Blue} 
        \textbf{Informazioni sul documento} & &\\ [1 ex]
        \rowcolor{LighterBlue}
        Destinatari: & Prof Tullio Vardanega & Prof Riccardo Cardin\\ [1 ex]
        \rowcolor{LightBlue}
        G al pedice: & Consultare il Glossario & \\ [1 ex]
    \end{tabular}
\end{center}
    
\newpage

\textbf{Registro dei Cambiamenti - Changelog\textsubscript{G}}
\begin{center}
\begin{tabular}{P{4em} P{6em} P{8em} P{8em} P{10em}} 
    \rowcolor{LightBlue}
    0.0.1 & 03/05/2024 & Giulia Dentone & Samuele Vignotto &  Definizione struttura del documento e scheletro delle sezioni. Scrittura introduzione ed obiettivi delle diverse sezioni\\
\end{tabular}
\end{center}
\newpage
\tableofcontents
\newpage

\section{Introduzione}
\subsection{Scopo del documento}
Scopo di questo documento è fornire agli sviluppatori le informazioni fondamentali per comprendere l'architettura del software\textsubscript{G} e le sue funzionalità. Inoltre, vengono illustrate le procedure per la configurazione, l'installazione e l'utilizzo del software\textsc{G}. Il Manuale Sviluppatore fornisce indicazioni sulla struttura del codice sorgente, sulle librerie e sulle API utilizzate, sui requisiti di sistema, sulla gestione degli errori e sui test da eseguire per garantire la qualità del software\textsubscript{G}. Rivolto principalmente a programmatori e sviluppatori che desiderano estendere o personalizzare il software\textsubscript{G}, il documento mira a fornire una comprensione approfondita della struttura e del funzionamento del prodotto. In tal modo, si assicura la corretta manutenzione e si forniscono dettagli essenziali per un successivo utilizzo e implementazione.
\subsection{Scopo del prodotto}
L'azienda proponente ha richiesto la creazione di una web app\textsubscript{G} che, tramite l'uso di IA\textsubscript{G} (in questo caso ChatGPT4 e Bedrock) è in grado di creare epic user stories\textsubscript{G} a partire dalle richieste del cliente e confrontarle con il codice sviluppato in modo da informare il cliente dello stato di avanzamento dello sviluppo del prodotto. Inoltre deve essere possibile, sia per il Project Manager\textsubscript{G}, sia per il cliente rilasciare dei feedback (nel primo caso riguardanti l'adeguatezza delle stories, nel secondo caso riguardanti il prodotto finale) al fine di migliorare l'IA\textsubscript{G}. È inoltre richiesta un' analisi comparativa tra le due IA\textsubscript{G} utilizzate e lo sviluppo di un plug-in\textsubscript{G} utile agli sviluppatori e al Project Manager\textsubscript{G}.

\subsection{Glossario}
Alcuni termini presenti nel documento potrebbero essere ambigui, pertanto verranno inseriti nel Glossario v.1.0.0. La loro presenza all'interno di esso sarà indicata tramite una G maiuscola a pedice. 

\section{Riferimenti}
\subsection{Riferimenti normativi}
\begin{itemize}
    \item Capitolato \textbf{C7 - ChatGPT vs BedRock developer Analysis}
    \\ \\
       \href{https://github.com/CyberSorceres/CyberSorceresRepository}{https://github.com/CyberSorceres/CyberSorceresRepository} 
    \item Norme del way of working v 1.0.0
    \item Regolamento del progetto didattico \\ \\ \href{https://www.math.unipd.it/~tullio/IS-1/2023/Dispense/PD2.pdf} 
    {https://www.math.unipd.it/~tullio/IS-1/2023/Dispense/PD2.pdf}
\end{itemize}
\subsection{Riferimenti informativi}
\begin{itemize}
    \item Slide del corso di Ingegneria del Software - Analisi dei requisiti \\ \\
    \href{https://www.math.unipd.it/~tullio/IS-1/2023/Dispense/T5.pdf}{https://www.math.unipd.it/~tullio/IS-1/2023/Dispense/T5.pdf}
    \item Slide del corso di Ingegneria del Software - Progettazione e programmazione: Diagrammi delle classi \\ \\
\href{https://www.math.unipd.it/~rcardin/swea/2023/Diagrammi%20delle%20Classi.pdf}{https://www.math.unipd.it/~rcardin/swea/2023/Diagrammi\%20delle\%20Classi.pdf}
    \item Slide del corso di Ingegneria del Software - Solid Programming \\ \\
\href{https://www.math.unipd.it/~rcardin/swea/2021/SOLID\%20Principles\%20of\%20Object-Oriented\%20Design_4x4.pdf}{\texttt{https://www.math.unipd.it/~rcardin/swea/2021/SOLID\%20Principles\%20of\%20Object-Oriented\%20Design\_4x4.pdf}}
\end{itemize}

\subsection{Riferimenti tecnici}
\begin{itemize}
\item Documentazione di React \\ \href{ https://react.dev/}{ https://react.dev/}
\item Documentazione di Typescript \\ \href{https://www.typescriptlang.org/docs/}{https://www.typescriptlang.org/docs/}
\item Documentazione di MongoDB \\ \href{https://www.mongodb.com/docs/}{https://www.mongodb.com/docs/}
\item Documentazione di Amazon AWS \\ \href{https://docs.aws.amazon.com/it_it/}{https://docs.aws.amazon.com/it\_it/}
\item Serverless Microservice Patterns \\\href{https://medium.com/@jeremydaly/serverless-microservice-patterns-for-aws-6dadcd21bc02}{https://medium.com/@jeremydaly/serverless-microservice-patterns-for-aws-6dadcd21bc02}
\item Aws Reference Architecture Diagrams \\ \href{https://aws.amazon.com/it/architecture/reference-architecture-diagrams}{https://aws.amazon.com/it/architecture/reference-architecture-diagrams}
\item React design patterns \\ \href{https://refine.dev/blog/react-design-patterns/}{https://refine.dev/blog/react-design-patterns/}
\end{itemize}

\section{Tecnologie}
In questa sezione viene presentata una visione d'insieme delle tecnologie adoperate per sviluppare il prodotto in questione. Vengono esaminate le tecnologie, gli strumenti e le librerie necessari per lo sviluppo, il testing\textsubscript{G} e la distribuzione del prodotto. Saranno in particolare affrontate le tecnologie impiegate per il front-end\textsubscript{G} e il back-end\textsubscript{G}, nonché per la gestione del database e l'integrazione con altri servizi.
\begin{center}
\begin{tabular}{c|c|c}
\hline
\rowcolor{Blue}
Tecnologia & Descrizione & Versione \\
\rowcolor{LighterBlue}

\rowcolor{LightBlue}
\end{tabular}
\end{center}
Inoltre, vengono definite le seguenti tecnologie per la fase di test\textsubscript{G}:
\begin{center}
\begin{tabular}{c|c|c}
\hline
\rowcolor{Blue}
Tecnologia & Descrizione & Versione \\
\rowcolor{LighterBlue}

\rowcolor{LightBlue}
\end{tabular}
\end{center}

\section{Configurazione}
Questa sezione fornisce un elenco dei requisiti minimi richiesti per eseguire l'applicazione, includendo le caratteristiche necessarie per configurare l'ambiente di sviluppo del progetto.
\subsection{Requisiti hardware}
\subsection{Requisiti software}

\section{Installazione}
\subsection{Clonazione della repository}
\subsection{Avvio dell'applicazione}

\section{Preparazione dell'ambiente di lavoro}
\section{Documentazione API}
\section{Punti di estensione dell'applicazione}


\end{document}
